\documentclass[11pt,a4paper]{article}
\usepackage{fontspec,amsmath,amssymb,amsthm,microtype}
\setmainfont{lmroman10-regular.otf}[BoldFont=lmroman10-bold.otf,ItalicFont=lmroman10-italic.otf,BoldItalicFont=lmroman10-bolditalic.otf]
\usepackage{xcolor}
\usepackage[most]{tcolorbox}
\usepackage{mathrsfs}
\providecommand{\lt}{<}
\providecommand{\gt}{>}
\setmonofont{DejaVuSansMono.ttf}[BoldFont=DejaVuSansMono-Bold.ttf,ItalicFont=DejaVuSansMono-Oblique.ttf,BoldItalicFont=DejaVuSansMono-BoldOblique.ttf,Scale=0.86]
\newfontfamily\leanfallback{STIXTwoMath-Regular.otf}[Scale=0.86]
\definecolor{leanrule}{gray}{0.62}
\definecolor{leanback}{gray}{0.962}
\definecolor{leantitle}{gray}{0.22}
\definecolor{leanmuted}{gray}{0.38}
\newtcolorbox{leanbox}[1]{enhanced,breakable,colback=leanback,colframe=leanrule,
  boxrule=0.4pt,arc=0.6mm,left=2.4mm,right=2.4mm,top=1.8mm,bottom=1.8mm,
  before skip=0.8em,after skip=1em,fonttitle=\small\bfseries,coltitle=leantitle,
  colbacktitle=leanback,titlerule=0.3pt,title={#1}}
\newlength{\leancharwidth}
\newenvironment{leancode}{\par\smallskip\begingroup\ttfamily\footnotesize\raggedright
  \settowidth{\leancharwidth}{M}\setlength{\parskip}{0pt}\setlength{\parindent}{0pt}
  \frenchspacing}{\par\endgroup\smallskip}
\newcommand{\leanline}[2]{\hangindent=\dimexpr#1\leancharwidth+4\leancharwidth\relax
  \hangafter=1\noindent\hspace*{#1\leancharwidth}#2\par}
\newcommand{\leansource}[1]{\par{\footnotesize\color{leanmuted}#1}\par}
\newcommand{\leanlabel}[1]{\par\smallskip\noindent{\small\bfseries #1}\par}
\usepackage[margin=25mm]{geometry}
\usepackage[unicode,hidelinks]{hyperref}
\hypersetup{pdfauthor={LangevinLean Project}}
\setlength{\emergencystretch}{3em}
\setlength{\parskip}{0.4em}
\setlength{\parindent}{0pt}
\allowdisplaybreaks[2]
\raggedbottom
\title{\texorpdfstring{Strong convexity forces \(m\le L\)}{Strong convexity forces m≤ L}}
\author{LangevinLean Project}
\date{}
\begin{document}\maketitle
Draft mathematical exposition; these notes may contain mistakes.


Let a differentiable potential on a nontrivial Hilbert space be \(m\)-strongly convex and have \(L\)-Lipschitz gradient. For distinct \(x,y\), strong monotonicity and Cauchy–Schwarz give \(m\|x-y\|^2\le\langle\nabla U(x)-\nabla U(y),x-y\rangle\le L\|x-y\|^2\). Dividing by the positive squared distance proves \(m\le L\). No positivity assumption on \(m\) is needed for this comparison.



Nontriviality is used to choose distinct points. In a singleton space every distance vanishes, and this argument cannot establish a relation between the two constants. The transport estimates on that space follow separately from zero coupling cost.


\section{\texorpdfstring{Compare monotonicity with Cauchy–Schwarz}{Compare monotonicity with Cauchy–Schwarz}}


Choose \(x\ne0\), possible because the Hilbert space is nontrivial. Strong monotonicity gives \(m\|x\|^2\le\langle\nabla U(x)-\nabla U(0),x\rangle\), while Cauchy–Schwarz and Lipschitzness bound the same quantity by \(L\|x\|^2\). Cancel the strictly positive \(\|x\|^2\) to prove \(m\le L\). No positivity premise on \(m\) is needed here.



The nontriviality assumption supplies a direction with positive squared norm. Dividing the two bounds for the same pairing by that norm gives the curvature comparison. On a singleton space no such direction exists, so the inequality need not follow; all transport errors there are instead treated directly as zero.

\begin{leanbox}{Lean: \texttt{strongConvex\_\allowbreak{}le\_\allowbreak{}lipschitz}}
\leansource{\texttt{strongConvex\_\allowbreak{}le\_\allowbreak{}lipschitz} \quad UnadjustedLangevin/\allowbreak{}PotentialAssumptions.lean, lines 25–40}
\begin{leancode}
\leanline{0}{theorem\ strongConvex\_le\_lipschitz\ \{U\ :\ H\ →\ ℝ\}\ \{m\ :\ ℝ\}\ \{L\ :\ ℝ≥0\}}
\leanline{4}{(hdiff\ :\ Differentiable\ ℝ\ U)\ (hsc\ :\ StrongConvexOn\ Set.univ\ m\ U)}
\leanline{4}{(hLip\ :\ LipschitzWith\ L\ (gradient\ U))\ :\ m\ ≤\ (L\ :\ ℝ)}
\end{leancode}
\leanlabel{How the Lean proof follows the argument}
\texttt{strong\allowbreak{}Convex\_\allowbreak{}le\_\allowbreak{}lipschitz}: the intermediate facts include \texttt{hxpos}, \texttt{hmono}, \texttt{hbound}, in source order; the source steps apply or rewrite \texttt{sub\_\allowbreak{}zero}, \texttt{strong\_\allowbreak{}gradient\_\allowbreak{}monotone}, \texttt{sq\_\allowbreak{}pos\_\allowbreak{}of\_\allowbreak{}pos}, \texttt{mul\_\allowbreak{}comm}, \texttt{hmono.\allowbreak{}trans}. The \texttt{calc} block checks the successive equalities or inequalities in the argument. The remaining scalar rearrangements are checked by \texttt{ring}.
\end{leanbox}

\end{document}
